% Options for packages loaded elsewhere
% Options for packages loaded elsewhere
\PassOptionsToPackage{unicode}{hyperref}
\PassOptionsToPackage{hyphens}{url}
\PassOptionsToPackage{dvipsnames,svgnames,x11names}{xcolor}
%
\documentclass[
  letterpaper,
  DIV=11,
  numbers=noendperiod]{scrartcl}
\usepackage{xcolor}
\usepackage{amsmath,amssymb}
\setcounter{secnumdepth}{5}
\usepackage{iftex}
\ifPDFTeX
  \usepackage[T1]{fontenc}
  \usepackage[utf8]{inputenc}
  \usepackage{textcomp} % provide euro and other symbols
\else % if luatex or xetex
  \usepackage{unicode-math} % this also loads fontspec
  \defaultfontfeatures{Scale=MatchLowercase}
  \defaultfontfeatures[\rmfamily]{Ligatures=TeX,Scale=1}
\fi
\usepackage{lmodern}
\ifPDFTeX\else
  % xetex/luatex font selection
\fi
% Use upquote if available, for straight quotes in verbatim environments
\IfFileExists{upquote.sty}{\usepackage{upquote}}{}
\IfFileExists{microtype.sty}{% use microtype if available
  \usepackage[]{microtype}
  \UseMicrotypeSet[protrusion]{basicmath} % disable protrusion for tt fonts
}{}
\makeatletter
\@ifundefined{KOMAClassName}{% if non-KOMA class
  \IfFileExists{parskip.sty}{%
    \usepackage{parskip}
  }{% else
    \setlength{\parindent}{0pt}
    \setlength{\parskip}{6pt plus 2pt minus 1pt}}
}{% if KOMA class
  \KOMAoptions{parskip=half}}
\makeatother
% Make \paragraph and \subparagraph free-standing
\makeatletter
\ifx\paragraph\undefined\else
  \let\oldparagraph\paragraph
  \renewcommand{\paragraph}{
    \@ifstar
      \xxxParagraphStar
      \xxxParagraphNoStar
  }
  \newcommand{\xxxParagraphStar}[1]{\oldparagraph*{#1}\mbox{}}
  \newcommand{\xxxParagraphNoStar}[1]{\oldparagraph{#1}\mbox{}}
\fi
\ifx\subparagraph\undefined\else
  \let\oldsubparagraph\subparagraph
  \renewcommand{\subparagraph}{
    \@ifstar
      \xxxSubParagraphStar
      \xxxSubParagraphNoStar
  }
  \newcommand{\xxxSubParagraphStar}[1]{\oldsubparagraph*{#1}\mbox{}}
  \newcommand{\xxxSubParagraphNoStar}[1]{\oldsubparagraph{#1}\mbox{}}
\fi
\makeatother

\usepackage{color}
\usepackage{fancyvrb}
\newcommand{\VerbBar}{|}
\newcommand{\VERB}{\Verb[commandchars=\\\{\}]}
\DefineVerbatimEnvironment{Highlighting}{Verbatim}{commandchars=\\\{\}}
% Add ',fontsize=\small' for more characters per line
\usepackage{framed}
\definecolor{shadecolor}{RGB}{241,243,245}
\newenvironment{Shaded}{\begin{snugshade}}{\end{snugshade}}
\newcommand{\AlertTok}[1]{\textcolor[rgb]{0.68,0.00,0.00}{#1}}
\newcommand{\AnnotationTok}[1]{\textcolor[rgb]{0.37,0.37,0.37}{#1}}
\newcommand{\AttributeTok}[1]{\textcolor[rgb]{0.40,0.45,0.13}{#1}}
\newcommand{\BaseNTok}[1]{\textcolor[rgb]{0.68,0.00,0.00}{#1}}
\newcommand{\BuiltInTok}[1]{\textcolor[rgb]{0.00,0.23,0.31}{#1}}
\newcommand{\CharTok}[1]{\textcolor[rgb]{0.13,0.47,0.30}{#1}}
\newcommand{\CommentTok}[1]{\textcolor[rgb]{0.37,0.37,0.37}{#1}}
\newcommand{\CommentVarTok}[1]{\textcolor[rgb]{0.37,0.37,0.37}{\textit{#1}}}
\newcommand{\ConstantTok}[1]{\textcolor[rgb]{0.56,0.35,0.01}{#1}}
\newcommand{\ControlFlowTok}[1]{\textcolor[rgb]{0.00,0.23,0.31}{\textbf{#1}}}
\newcommand{\DataTypeTok}[1]{\textcolor[rgb]{0.68,0.00,0.00}{#1}}
\newcommand{\DecValTok}[1]{\textcolor[rgb]{0.68,0.00,0.00}{#1}}
\newcommand{\DocumentationTok}[1]{\textcolor[rgb]{0.37,0.37,0.37}{\textit{#1}}}
\newcommand{\ErrorTok}[1]{\textcolor[rgb]{0.68,0.00,0.00}{#1}}
\newcommand{\ExtensionTok}[1]{\textcolor[rgb]{0.00,0.23,0.31}{#1}}
\newcommand{\FloatTok}[1]{\textcolor[rgb]{0.68,0.00,0.00}{#1}}
\newcommand{\FunctionTok}[1]{\textcolor[rgb]{0.28,0.35,0.67}{#1}}
\newcommand{\ImportTok}[1]{\textcolor[rgb]{0.00,0.46,0.62}{#1}}
\newcommand{\InformationTok}[1]{\textcolor[rgb]{0.37,0.37,0.37}{#1}}
\newcommand{\KeywordTok}[1]{\textcolor[rgb]{0.00,0.23,0.31}{\textbf{#1}}}
\newcommand{\NormalTok}[1]{\textcolor[rgb]{0.00,0.23,0.31}{#1}}
\newcommand{\OperatorTok}[1]{\textcolor[rgb]{0.37,0.37,0.37}{#1}}
\newcommand{\OtherTok}[1]{\textcolor[rgb]{0.00,0.23,0.31}{#1}}
\newcommand{\PreprocessorTok}[1]{\textcolor[rgb]{0.68,0.00,0.00}{#1}}
\newcommand{\RegionMarkerTok}[1]{\textcolor[rgb]{0.00,0.23,0.31}{#1}}
\newcommand{\SpecialCharTok}[1]{\textcolor[rgb]{0.37,0.37,0.37}{#1}}
\newcommand{\SpecialStringTok}[1]{\textcolor[rgb]{0.13,0.47,0.30}{#1}}
\newcommand{\StringTok}[1]{\textcolor[rgb]{0.13,0.47,0.30}{#1}}
\newcommand{\VariableTok}[1]{\textcolor[rgb]{0.07,0.07,0.07}{#1}}
\newcommand{\VerbatimStringTok}[1]{\textcolor[rgb]{0.13,0.47,0.30}{#1}}
\newcommand{\WarningTok}[1]{\textcolor[rgb]{0.37,0.37,0.37}{\textit{#1}}}

\usepackage{longtable,booktabs,array}
\newcounter{none} % for unnumbered tables
\usepackage{calc} % for calculating minipage widths
% Correct order of tables after \paragraph or \subparagraph
\usepackage{etoolbox}
\makeatletter
\patchcmd\longtable{\par}{\if@noskipsec\mbox{}\fi\par}{}{}
\makeatother
% Allow footnotes in longtable head/foot
\IfFileExists{footnotehyper.sty}{\usepackage{footnotehyper}}{\usepackage{footnote}}
\makesavenoteenv{longtable}
\usepackage{graphicx}
\makeatletter
\newsavebox\pandoc@box
\newcommand*\pandocbounded[1]{% scales image to fit in text height/width
  \sbox\pandoc@box{#1}%
  \Gscale@div\@tempa{\textheight}{\dimexpr\ht\pandoc@box+\dp\pandoc@box\relax}%
  \Gscale@div\@tempb{\linewidth}{\wd\pandoc@box}%
  \ifdim\@tempb\p@<\@tempa\p@\let\@tempa\@tempb\fi% select the smaller of both
  \ifdim\@tempa\p@<\p@\scalebox{\@tempa}{\usebox\pandoc@box}%
  \else\usebox{\pandoc@box}%
  \fi%
}
% Set default figure placement to htbp
\def\fps@figure{htbp}
\makeatother


% definitions for citeproc citations
\NewDocumentCommand\citeproctext{}{}
\NewDocumentCommand\citeproc{mm}{%
  \begingroup\def\citeproctext{#2}\cite{#1}\endgroup}
\makeatletter
 % allow citations to break across lines
 \let\@cite@ofmt\@firstofone
 % avoid brackets around text for \cite:
 \def\@biblabel#1{}
 \def\@cite#1#2{{#1\if@tempswa , #2\fi}}
\makeatother
\newlength{\cslhangindent}
\setlength{\cslhangindent}{1.5em}
\newlength{\csllabelwidth}
\setlength{\csllabelwidth}{3em}
\newenvironment{CSLReferences}[2] % #1 hanging-indent, #2 entry-spacing
 {\begin{list}{}{%
  \setlength{\itemindent}{0pt}
  \setlength{\leftmargin}{0pt}
  \setlength{\parsep}{0pt}
  % turn on hanging indent if param 1 is 1
  \ifodd #1
   \setlength{\leftmargin}{\cslhangindent}
   \setlength{\itemindent}{-1\cslhangindent}
  \fi
  % set entry spacing
  \setlength{\itemsep}{#2\baselineskip}}}
 {\end{list}}
\usepackage{calc}
\newcommand{\CSLBlock}[1]{\hfill\break\parbox[t]{\linewidth}{\strut\ignorespaces#1\strut}}
\newcommand{\CSLLeftMargin}[1]{\parbox[t]{\csllabelwidth}{\strut#1\strut}}
\newcommand{\CSLRightInline}[1]{\parbox[t]{\linewidth - \csllabelwidth}{\strut#1\strut}}
\newcommand{\CSLIndent}[1]{\hspace{\cslhangindent}#1}



\setlength{\emergencystretch}{3em} % prevent overfull lines

\providecommand{\tightlist}{%
  \setlength{\itemsep}{0pt}\setlength{\parskip}{0pt}}



 


\KOMAoption{captions}{tableheading}
\usepackage{fontspec}
\newfontfamily\ipaFont{Charis}
\newcommand{\ipa}[1]{{\ipaFont #1}}
\makeatletter
\renewcommand{\tableofcontents}{%
  \section*{Table of Contents}%
  \@starttoc{toc}%
}
\makeatother
\makeatletter
\@ifpackageloaded{caption}{}{\usepackage{caption}}
\AtBeginDocument{%
\ifdefined\contentsname
  \renewcommand*\contentsname{Table of contents}
\else
  \newcommand\contentsname{Table of contents}
\fi
\ifdefined\listfigurename
  \renewcommand*\listfigurename{List of Figures}
\else
  \newcommand\listfigurename{List of Figures}
\fi
\ifdefined\listtablename
  \renewcommand*\listtablename{List of Tables}
\else
  \newcommand\listtablename{List of Tables}
\fi
\ifdefined\figurename
  \renewcommand*\figurename{Figure}
\else
  \newcommand\figurename{Figure}
\fi
\ifdefined\tablename
  \renewcommand*\tablename{Table}
\else
  \newcommand\tablename{Table}
\fi
}
\@ifpackageloaded{float}{}{\usepackage{float}}
\floatstyle{ruled}
\@ifundefined{c@chapter}{\newfloat{codelisting}{h}{lop}}{\newfloat{codelisting}{h}{lop}[chapter]}
\floatname{codelisting}{Listing}
\newcommand*\listoflistings{\listof{codelisting}{List of Listings}}
\makeatother
\makeatletter
\makeatother
\makeatletter
\@ifpackageloaded{caption}{}{\usepackage{caption}}
\@ifpackageloaded{subcaption}{}{\usepackage{subcaption}}
\makeatother
\usepackage{bookmark}
\IfFileExists{xurl.sty}{\usepackage{xurl}}{} % add URL line breaks if available
\urlstyle{same}
\hypersetup{
  pdftitle={Preprocessing\_data},
  pdfauthor={Tiziana Ilie},
  colorlinks=true,
  linkcolor={blue},
  filecolor={Maroon},
  citecolor={Blue},
  urlcolor={Blue},
  pdfcreator={LaTeX via pandoc}}


\title{Preprocessing\_data}
\author{Tiziana Ilie}
\date{}
\begin{document}
\maketitle


\section{Set-up}\label{set-up}

\begin{Shaded}
\begin{Highlighting}[]
\CommentTok{\#loading packages}

\FunctionTok{library}\NormalTok{(here)}
\FunctionTok{library}\NormalTok{(tidyverse)}
\end{Highlighting}
\end{Shaded}

\section{Preprocessing the Spanish native
data}\label{preprocessing-the-spanish-native-data}

\begin{Shaded}
\begin{Highlighting}[]
\CommentTok{\#Loading Kocher\textquotesingle{}s Spanish native data}

\NormalTok{kocher\_native\_data }\OtherTok{\textless{}{-}} \FunctionTok{read.csv}\NormalTok{(here}\SpecialCharTok{::}\FunctionTok{here}\NormalTok{(}\StringTok{"data"}\NormalTok{, }\StringTok{"raw\_data"}\NormalTok{, }\StringTok{"datanatives.csv"}\NormalTok{))}

\CommentTok{\# Adding the \textquotesingle{}SubjectGenus\textquotesingle{} variable to the Kocher\textquotesingle{}s dataframe}

\NormalTok{kocher\_native\_data[, }\StringTok{\textquotesingle{}SubjectGenus\textquotesingle{}}\NormalTok{] }\OtherTok{=} \ConstantTok{NA}

\CommentTok{\#write\_csv(kocher\_native\_data, here("data", "kocher\_native\_data.csv")) \#saving added data as .csv file}
\end{Highlighting}
\end{Shaded}

\section{Preprocessing the L1 German
data}\label{preprocessing-the-l1-german-data}

\begin{Shaded}
\begin{Highlighting}[]
\CommentTok{\#Loading Kocher\textquotesingle{}s L1 German data and the additionally extracted L1 German data with more speaker texts}

\NormalTok{raw\_data\_germanL1 }\OtherTok{\textless{}{-}} \FunctionTok{read.csv}\NormalTok{(here}\SpecialCharTok{::}\FunctionTok{here}\NormalTok{(}\StringTok{"data"}\NormalTok{, }\StringTok{"raw\_data"}\NormalTok{, }\StringTok{"raw\_data\_germanL1.csv"}\NormalTok{),}\AttributeTok{sep =} \StringTok{"}\SpecialCharTok{\textbackslash{}t}\StringTok{"}\NormalTok{) }\CommentTok{\#extracted data from CEDEL2 at 2026{-}03{-}31}

\NormalTok{kocher\_data\_germanL1 }\OtherTok{\textless{}{-}} \FunctionTok{read.csv}\NormalTok{(here}\SpecialCharTok{::}\FunctionTok{here}\NormalTok{(}\StringTok{"data"}\NormalTok{, }\StringTok{"kocher\_germanL1\_data.csv"}\NormalTok{)) }\CommentTok{\#Kocher\textquotesingle{}s annotated data}

\CommentTok{\#Extracting the newly added speakers}
\NormalTok{new\_germanL1 }\OtherTok{\textless{}{-}} \FunctionTok{anti\_join}\NormalTok{(raw\_data\_germanL1, kocher\_data\_germanL1, }
          \AttributeTok{by =} \StringTok{"Filename"}\NormalTok{) }\CommentTok{\#17 in total, not segmented nor annotated yet}
\end{Highlighting}
\end{Shaded}

The next step consists of some data wrangling:

\begin{Shaded}
\begin{Highlighting}[]
\CommentTok{\# Changing abbreviations of column \textquotesingle{}Proficiency\textquotesingle{} to full speaker proficiency level names in the dataframe \textquotesingle{}germanL1\_data\_full\textquotesingle{} and adding the column \textquotesingle{}ProficiencyNr\textquotesingle{} to better understand and to later merch the data of \textquotesingle{}new\_germanL1\textquotesingle{} into the dataframe \textquotesingle{}kocher\_data\_germanL1\textquotesingle{}. }

\NormalTok{kocher\_data\_germanL1 }\OtherTok{\textless{}{-}}\NormalTok{ kocher\_data\_germanL1 }\SpecialCharTok{|\textgreater{}} 
  \FunctionTok{mutate}\NormalTok{(}
    \AttributeTok{Proficiency =} \FunctionTok{case\_when}\NormalTok{(}
\NormalTok{      Proficiency }\SpecialCharTok{==} \StringTok{"LB"} \SpecialCharTok{\textasciitilde{}} \StringTok{"Lower beginner"}\NormalTok{,}
\NormalTok{      Proficiency }\SpecialCharTok{==} \StringTok{"UB"} \SpecialCharTok{\textasciitilde{}} \StringTok{"Upper beginner"}\NormalTok{,}
\NormalTok{      Proficiency }\SpecialCharTok{==} \StringTok{"LI"} \SpecialCharTok{\textasciitilde{}} \StringTok{"Lower intermediate"}\NormalTok{,}
\NormalTok{      Proficiency }\SpecialCharTok{==} \StringTok{"UI"} \SpecialCharTok{\textasciitilde{}} \StringTok{"Upper intermediate"}\NormalTok{,}
\NormalTok{      Proficiency }\SpecialCharTok{==} \StringTok{"LA"} \SpecialCharTok{\textasciitilde{}} \StringTok{"Lower advanced"}\NormalTok{,}
\NormalTok{      Proficiency }\SpecialCharTok{==} \StringTok{"UA"} \SpecialCharTok{\textasciitilde{}} \StringTok{"Upper advanced"}\NormalTok{,}
      \ConstantTok{TRUE} \SpecialCharTok{\textasciitilde{}}\NormalTok{ Proficiency)}
\NormalTok{  )}
\end{Highlighting}
\end{Shaded}

Further, I need to change the column names of the dataframe
`new\_germanL1' so that they agree with those of the dataframe
`kocher\_data\_germanL1' and deselect irrelevant columns in
`new\_germanL1' df. After that, I add the language proficiency level in
numbers (level 1-6), in addition to the different levels names.

\begin{Shaded}
\begin{Highlighting}[]
\FunctionTok{colnames}\NormalTok{(new\_germanL1) }\CommentTok{\#current column names}
\end{Highlighting}
\end{Shaded}

\begin{verbatim}
 [1] "Subcorpus"                                                     
 [2] "Filename"                                                      
 [3] "Year.data.collection"                                          
 [4] "Placement.test.score..raw."                                    
 [5] "Placement.test.score...."                                      
 [6] "Proficiency"                                                   
 [7] "Sex"                                                           
 [8] "Age"                                                           
 [9] "School.University.Institution"                                 
[10] "Major"                                                         
[11] "Year.at.university.school"                                     
[12] "L1"                                                            
[13] "Father.s.native.language"                                      
[14] "Mother.s.nativelanguage"                                       
[15] "Languages.spoken.at.home"                                      
[16] "Age.of.exposure.to.Spanish"                                    
[17] "Years.studying.Spanish"                                        
[18] "Stay.abroad.in.Spanish.speaking.country.....1.month."          
[19] "Stay.abroad..where."                                           
[20] "Stay.abroad..when."                                            
[21] "Stay.abroad..months."                                          
[22] "Language.certificates..type.and.level."                        
[23] "Proficiency..self.assessment..speaking"                        
[24] "Proficiency..self.assessment..listening"                       
[25] "Proficiency..self.assessment..reading"                         
[26] "Proficiency..self.assessment..writing"                         
[27] "Proficiency..self.assessment."                                 
[28] "Additional.foreign.language.s."                                
[29] "Proficiency..self.assessment..in.additional.language.speaking" 
[30] "Proficiency..self.assessment..in.additional.language.listening"
[31] "Proficiency..self.assessment..in.additional.language.reading"  
[32] "Proficiency..self.assessment..in.additional.language.writing"  
[33] "Medium"                                                        
[34] "Task.number"                                                   
[35] "Task.title"                                                    
[36] "Writting.audio.details"                                        
[37] "Minutes.taken.to.complete.the.task"                            
[38] "Where.the.task.was.done"                                       
[39] "Resources.used"                                                
[40] "Text"                                                          
[41] "Original.text"                                                 
\end{verbatim}

\begin{Shaded}
\begin{Highlighting}[]
\NormalTok{new\_germanL1 }\OtherTok{\textless{}{-}}\NormalTok{ new\_germanL1 }\SpecialCharTok{|\textgreater{}}  \CommentTok{\#renaming columns as in Kochers df}
  \FunctionTok{rename}\NormalTok{(}
    \StringTok{"Year"} \OtherTok{=} \StringTok{"Year.data.collection"}\NormalTok{,}
    \StringTok{"AgeExposureSpanish"} \OtherTok{=} \StringTok{"Age.of.exposure.to.Spanish"}\NormalTok{,}
    \StringTok{"YearsstudySpanish"} \OtherTok{=} \StringTok{"Years.studying.Spanish"}\NormalTok{,}
    \StringTok{"AbroadSpanishCountry"} \OtherTok{=} \StringTok{"Stay.abroad.in.Spanish.speaking.country.....1.month."}\NormalTok{,}
    \StringTok{"DurationTask"} \OtherTok{=} \StringTok{"Minutes.taken.to.complete.the.task"}\NormalTok{,}
    \StringTok{"LocationTask"} \OtherTok{=} \StringTok{"Where.the.task.was.done"}\NormalTok{,}
    \StringTok{"ForeignLanguage"} \OtherTok{=} \StringTok{"Additional.foreign.language.s."}\NormalTok{,}
    \StringTok{"Score"} \OtherTok{=} \StringTok{"Placement.test.score...."}\NormalTok{,}
    \StringTok{"Place"} \OtherTok{=} \StringTok{"School.University.Institution"}\NormalTok{,}
    \StringTok{"FullText"} \OtherTok{=} \StringTok{"Text"}
\NormalTok{  ) }\SpecialCharTok{|\textgreater{}} 
  \FunctionTok{select}\NormalTok{( }\CommentTok{\#deselecting columns not needed}
    \SpecialCharTok{{-}}\NormalTok{Subcorpus, }\SpecialCharTok{{-}}\NormalTok{Placement.test.score..raw., }\SpecialCharTok{{-}}\NormalTok{L1, }\SpecialCharTok{{-}}\NormalTok{Major, }\SpecialCharTok{{-}}\NormalTok{Year.at.university.school, }\SpecialCharTok{{-}}\NormalTok{Father.s.native.language, }\SpecialCharTok{{-}}\NormalTok{Mother.s.nativelanguage, }\SpecialCharTok{{-}}\NormalTok{Stay.abroad..where., }\SpecialCharTok{{-}}\NormalTok{Stay.abroad..when., }\SpecialCharTok{{-}}\NormalTok{Stay.abroad..months., }\SpecialCharTok{{-}}\NormalTok{Language.certificates..type.and.level., }\SpecialCharTok{{-}}\NormalTok{Proficiency..self.assessment..speaking, }\SpecialCharTok{{-}}\NormalTok{Proficiency..self.assessment..listening, }\SpecialCharTok{{-}}\NormalTok{Proficiency..self.assessment..reading, }\SpecialCharTok{{-}}\NormalTok{Proficiency..self.assessment..writing, }\SpecialCharTok{{-}}\NormalTok{Proficiency..self.assessment., }\SpecialCharTok{{-}}\NormalTok{Proficiency..self.assessment..in.additional.language.speaking, }\SpecialCharTok{{-}}\NormalTok{Proficiency..self.assessment..in.additional.language.listening, }\SpecialCharTok{{-}}\NormalTok{Proficiency..self.assessment..in.additional.language.reading, }\SpecialCharTok{{-}}\NormalTok{Proficiency..self.assessment..in.additional.language.writing, }\SpecialCharTok{{-}}\NormalTok{Medium, }\SpecialCharTok{{-}}\NormalTok{Task.number, }\SpecialCharTok{{-}}\NormalTok{Task.title, }\SpecialCharTok{{-}}\NormalTok{Writting.audio.details, }\SpecialCharTok{{-}}\NormalTok{Resources.used, }\SpecialCharTok{{-}}\NormalTok{Original.text) }\SpecialCharTok{|\textgreater{}} 
  \FunctionTok{mutate}\NormalTok{( }\CommentTok{\#adding ProficiencyNr column and values 1{-}6}
    \AttributeTok{ProficiencyNr =} \FunctionTok{case\_when}\NormalTok{(}
\NormalTok{      Proficiency }\SpecialCharTok{==} \StringTok{"Lower beginner"} \SpecialCharTok{\textasciitilde{}} \StringTok{"1"}\NormalTok{,}
\NormalTok{      Proficiency }\SpecialCharTok{==} \StringTok{"Upper beginner"} \SpecialCharTok{\textasciitilde{}} \StringTok{"2"}\NormalTok{,}
\NormalTok{      Proficiency }\SpecialCharTok{==} \StringTok{"Lower intermediate"} \SpecialCharTok{\textasciitilde{}} \StringTok{"3"}\NormalTok{,}
\NormalTok{      Proficiency }\SpecialCharTok{==} \StringTok{"Upper intermediate"}\SpecialCharTok{\textasciitilde{}} \StringTok{"4"}\NormalTok{,}
\NormalTok{      Proficiency }\SpecialCharTok{==} \StringTok{"Lower advanced"} \SpecialCharTok{\textasciitilde{}} \StringTok{"5"}\NormalTok{,}
\NormalTok{      Proficiency }\SpecialCharTok{==} \StringTok{"Upper advanced"} \SpecialCharTok{\textasciitilde{}} \StringTok{"6"}\NormalTok{,}
      \ConstantTok{TRUE} \SpecialCharTok{\textasciitilde{}}\NormalTok{ Proficiency)}
\NormalTok{    )}
\end{Highlighting}
\end{Shaded}

Next, I want to add the `NSLhome' and `NSLadditional' columns and
therefore extract the information of the column
`Languages.spoken.at.home'.

\begin{Shaded}
\begin{Highlighting}[]
\FunctionTok{table}\NormalTok{(new\_germanL1}\SpecialCharTok{$}\NormalTok{Languages.spoken.at.home) }\CommentTok{\# displaying languages spoken at home}
\end{Highlighting}
\end{Shaded}

\begin{verbatim}

          German  German, English  German, Italian  German, Serbian 
              10                1                1                1 
         Spanish Spanish, English  Spanish, German Turkish, Spanish 
               1                1                1                1 
\end{verbatim}

\begin{Shaded}
\begin{Highlighting}[]
\CommentTok{\#Apart from German, some speakers also grew up acquiring English, Italian, Serbian, Spanish and Turkish}

\FunctionTok{class}\NormalTok{(new\_germanL1}\SpecialCharTok{$}\NormalTok{Languages.spoken.at.home) }\CommentTok{\#vector type "character"}
\end{Highlighting}
\end{Shaded}

\begin{verbatim}
[1] "character"
\end{verbatim}

\begin{Shaded}
\begin{Highlighting}[]
\NormalTok{nsl\_home\_L1germ }\OtherTok{\textless{}{-}} \FunctionTok{c}\NormalTok{(}\StringTok{"Spanish"}\NormalTok{, }\StringTok{"Italian"}\NormalTok{, }\StringTok{"Serbian"}\NormalTok{, }\StringTok{"Turkish"}\NormalTok{)}

\NormalTok{new\_germanL1 }\OtherTok{\textless{}{-}}\NormalTok{ new\_germanL1 }\SpecialCharTok{|\textgreater{}}
  \FunctionTok{mutate}\NormalTok{(}
    \AttributeTok{NSLhome =} \FunctionTok{case\_when}\NormalTok{(}
      \FunctionTok{str\_detect}\NormalTok{(Languages.spoken.at.home, }\FunctionTok{paste}\NormalTok{(nsl\_home\_L1germ, }\AttributeTok{collapse =} \StringTok{"|"}\NormalTok{)) }\SpecialCharTok{\textasciitilde{}} \StringTok{"Yes"}\NormalTok{,}
      \ConstantTok{TRUE} \SpecialCharTok{\textasciitilde{}} \StringTok{"No"}
\NormalTok{    )}
\NormalTok{  ) }\CommentTok{\#If the speakers grew up with a NSL at home the answer is "Yes", else "No"}

\CommentTok{\#Adding the NSLadditional column}

\FunctionTok{class}\NormalTok{(new\_germanL1}\SpecialCharTok{$}\NormalTok{ForeignLanguage) }\CommentTok{\#vector type "character"}
\end{Highlighting}
\end{Shaded}

\begin{verbatim}
[1] "character"
\end{verbatim}

\begin{Shaded}
\begin{Highlighting}[]
\FunctionTok{table}\NormalTok{(new\_germanL1}\SpecialCharTok{$}\NormalTok{ForeignLanguage) }\CommentTok{\# displaying languages spoken at home}
\end{Highlighting}
\end{Shaded}

\begin{verbatim}

English Italian 
     16       1 
\end{verbatim}

\begin{Shaded}
\begin{Highlighting}[]
\NormalTok{new\_germanL1 }\OtherTok{\textless{}{-}}\NormalTok{ new\_germanL1 }\SpecialCharTok{|\textgreater{}}
  \FunctionTok{mutate}\NormalTok{(}\AttributeTok{NSLadditional =} \FunctionTok{case\_when}\NormalTok{(}
\NormalTok{    NSLhome }\SpecialCharTok{==} \StringTok{"Yes"} \SpecialCharTok{\textasciitilde{}} \StringTok{"Yes"}\NormalTok{,}
\NormalTok{    NSLhome }\SpecialCharTok{==} \StringTok{"No"} \SpecialCharTok{\&}\NormalTok{ ForeignLanguage }\SpecialCharTok{==} \StringTok{"Italien"} \SpecialCharTok{\textasciitilde{}} \StringTok{"Yes"}\NormalTok{,}
    \ConstantTok{TRUE} \SpecialCharTok{\textasciitilde{}} \StringTok{"No"} 
\NormalTok{  )) }\CommentTok{\#When speakers already acquired a NSL at home than column "NSLadditional" will already be "yes". If they did not but learned a NSL later (here only one speaker learned Italian, the others only English which is not a NSL), than in column "NSLadditional" will be annotated "yes".}
\end{Highlighting}
\end{Shaded}

Before merging the two dataframes I need the values in the column
`ProficiencyNr' to be the same vector type:

\begin{Shaded}
\begin{Highlighting}[]
\FunctionTok{class}\NormalTok{(kocher\_data\_germanL1}\SpecialCharTok{$}\NormalTok{ProficiencyNr) }\CommentTok{\#integer}
\end{Highlighting}
\end{Shaded}

\begin{verbatim}
[1] "integer"
\end{verbatim}

\begin{Shaded}
\begin{Highlighting}[]
\FunctionTok{class}\NormalTok{(new\_germanL1}\SpecialCharTok{$}\NormalTok{ProficiencyNr) }\CommentTok{\#character}
\end{Highlighting}
\end{Shaded}

\begin{verbatim}
[1] "character"
\end{verbatim}

\begin{Shaded}
\begin{Highlighting}[]
\NormalTok{new\_germanL1 }\OtherTok{\textless{}{-}}\NormalTok{ new\_germanL1 }\SpecialCharTok{|\textgreater{}} 
  \FunctionTok{mutate}\NormalTok{(}\AttributeTok{ProficiencyNr =} \FunctionTok{as.integer}\NormalTok{(ProficiencyNr)) }\CommentTok{\#integer now}

\CommentTok{\#write\_csv(new\_germanL1, here("data", "new\_germanL1.csv")) \#saving added data as .csv file}
\end{Highlighting}
\end{Shaded}

Now I can segment each clause per row of every text. Therefore I use the
Large Language Models (LLMs) \emph{Meta Llama 3.3 70B} and
\emph{DeepSeek R1 Distill LLama 70B} Instruct for this process via the
University of Cologne's academic cloud access. As this copus contains
learner data most of the clauses were not segmented correctly which is
why I had to manually recognise and segment the remaining clauses. The
new\_germanL1\_segmented dataframe is then read into R again for
merging.

\begin{Shaded}
\begin{Highlighting}[]
\CommentTok{\#Reading the manually segmented file in R}

\NormalTok{new\_germanL1\_segmented }\OtherTok{\textless{}{-}} \FunctionTok{read.csv}\NormalTok{(here}\SpecialCharTok{::}\FunctionTok{here}\NormalTok{(}\StringTok{"data"}\NormalTok{,}\StringTok{"new\_germanL1\_segmented.csv"}\NormalTok{), }\AttributeTok{stringsAsFactors=}\ConstantTok{TRUE}\NormalTok{)}

\CommentTok{\# Combining the segmented clauses of the dataframe \textquotesingle{}new\_germanL1\_segmented\textquotesingle{} with the metadata of the dataframe \textquotesingle{}new\_germanL1\textquotesingle{}.}

\NormalTok{joined\_new\_germanL1\_segmented }\OtherTok{\textless{}{-}} \FunctionTok{inner\_join}\NormalTok{(}\AttributeTok{x =}\NormalTok{ new\_germanL1\_segmented, }\AttributeTok{y =}\NormalTok{ new\_germanL1, }\AttributeTok{by =} \StringTok{"Filename"}\NormalTok{)}

\CommentTok{\#I obtain a dataframe with 18 variables/columns as \textquotesingle{}FullText\textquotesingle{} was added to the remaining 17 variables/columns of the dataframe \textquotesingle{}new\_germanL1\textquotesingle{}.}
\end{Highlighting}
\end{Shaded}

Now, I merge the dataframe `joined\_new\_germanL1\_segmented' containing
only the clauses, filenames and metadata of the new added L1German data
to Kocher (\citeproc{ref-kocher2025}{2025b})'s already annotated clauses
using the function bind\_rows(). The new dataframe's name is
\emph{joined\_full\_germanL1}.

\begin{Shaded}
\begin{Highlighting}[]
\NormalTok{joined\_germanL1\_data }\OtherTok{\textless{}{-}} \FunctionTok{bind\_rows}\NormalTok{(kocher\_data\_germanL1, joined\_new\_germanL1\_segmented)}

\CommentTok{\#Afterwards the column \textquotesingle{}Thetarole\textquotesingle{} was deselected as I consider this to be a very subjective variable and it also did not enter into Kocher\textquotesingle{}s model. Further, the column \textquotesingle{}SentenceType\textquotesingle{} was also deleted as the majority of the entries in this column are declarative sentences. }

\NormalTok{joined\_germanL1\_data }\OtherTok{\textless{}{-}}\NormalTok{ joined\_germanL1\_data }\SpecialCharTok{|\textgreater{}} 
  \FunctionTok{select}\NormalTok{(}\SpecialCharTok{{-}}\NormalTok{Thetarole,}\SpecialCharTok{{-}}\NormalTok{SentenceType) }

\CommentTok{\#Adding the \textquotesingle{}SubjectGenus\textquotesingle{} variable to the \textquotesingle{}joined\_germanL1\_data\textquotesingle{} dataset}

\NormalTok{joined\_germanL1\_data[, }\StringTok{\textquotesingle{}SubjectGenus\textquotesingle{}}\NormalTok{] }\OtherTok{=} \ConstantTok{NA}

\CommentTok{\#write\_csv(joined\_germanL1\_data, here("data", "joined\_germanL1\_data.csv")) \#saving added data as .csv file}
\end{Highlighting}
\end{Shaded}

This .csv file can now be manually annotated. Concretely, I annotate the
values of the columns:

\begin{itemize}
\item
  `Textnr'
\item
  `Sentence'
\item
  `ID'
\item
  `SubjectPosition'
\item
  `ClauseType'
\item
  `Specificity'
\item
  `InformationStructure'
\item
  `ReferentialContinuity'
\item
  `Subjecttype'
\item
  `Verb'
\item
  `VerbConstruction'
\item
  `SubjectGenus'
\end{itemize}

\section{L1 German data extraction for checking on the robustness of
Kocher's
annotation}\label{l1-german-data-extraction-for-checking-on-the-robustness-of-kochers-annotation}

For a later robustness check I will randomly extract 3 speakers of
Kocher (\citeproc{ref-kocher2025}{2025b})'s dataset and re-annotate
their clauses in order to check if the annotation processes align.

\begin{Shaded}
\begin{Highlighting}[]
\CommentTok{\#Letting R select three Filenames randomly}

\FunctionTok{set.seed}\NormalTok{(}\DecValTok{8}\NormalTok{)  }\CommentTok{\#Setting a seed to ensure that the outcome of my randomly generated number series are always exactly the same, so that my study is reproducible }
\NormalTok{random\_filenames }\OtherTok{\textless{}{-}} \FunctionTok{sample}\NormalTok{(kocher\_data\_germanL1}\SpecialCharTok{$}\NormalTok{Filename, }\DecValTok{3}\NormalTok{)}

\CommentTok{\#[1] "DE\_WR\_41\_23\_5\_14\_PN" }
\CommentTok{\#[2] "DE\_WR\_39\_46\_30\_14\_IK"}
\CommentTok{\#[3] "DE\_WR\_41\_21\_3\_14\_MP" }


\CommentTok{\#Creating a new dataframe with only the rows of the three selected filenames}

\NormalTok{sample\_kocher\_germanL1 }\OtherTok{\textless{}{-}}\NormalTok{ kocher\_data\_germanL1 }\SpecialCharTok{|\textgreater{}} 
  \FunctionTok{filter}\NormalTok{(Filename }\SpecialCharTok{\%in\%} \FunctionTok{c}\NormalTok{(}
    \StringTok{"DE\_WR\_41\_23\_5\_14\_PN"}\NormalTok{,}
    \StringTok{"DE\_WR\_39\_46\_30\_14\_IK"}\NormalTok{,}
    \StringTok{"DE\_WR\_41\_21\_3\_14\_MP"}\NormalTok{)}
\NormalTok{   )}

\CommentTok{\# Adding the \textquotesingle{}SubjectGenus\textquotesingle{} column}

\NormalTok{sample\_kocher\_germanL1[, }\StringTok{\textquotesingle{}SubjectGenus\textquotesingle{}}\NormalTok{] }\OtherTok{=} \ConstantTok{NA}

\CommentTok{\# Deleting \textquotesingle{}Thetarole\textquotesingle{} and \textquotesingle{}SentenceType\textquotesingle{} columns}

\NormalTok{sample\_kocher\_germanL1 }\OtherTok{\textless{}{-}}\NormalTok{ sample\_kocher\_germanL1 }\SpecialCharTok{|\textgreater{}} 
  \FunctionTok{select}\NormalTok{(}\SpecialCharTok{{-}}\NormalTok{Thetarole, }\SpecialCharTok{{-}}\NormalTok{SentenceType) }\CommentTok{\#deselecting columns not needed}

\CommentTok{\#write\_csv(sample\_kocher\_germanL1, here("data", "sample\_kocher\_germanL1.csv")) \#saving data as .csv file}
\end{Highlighting}
\end{Shaded}

The columns of df ``sample\_kocher\_germanL1'' will be emptied and
reannotated manually.

\section{Preprocessing the L1 English
data}\label{preprocessing-the-l1-english-data}

\begin{Shaded}
\begin{Highlighting}[]
\NormalTok{englishL1\_data }\OtherTok{\textless{}{-}} \FunctionTok{read.csv}\NormalTok{(here}\SpecialCharTok{::}\FunctionTok{here}\NormalTok{(}\StringTok{"data"}\NormalTok{, }\StringTok{"raw\_data/raw\_data\_englishL1.csv"}\NormalTok{), }\AttributeTok{sep =} \StringTok{"}\SpecialCharTok{\textbackslash{}t}\StringTok{"}\NormalTok{)}

\CommentTok{\# Changing the column names so they align with the L1German data.}

\FunctionTok{colnames}\NormalTok{(englishL1\_data) }\CommentTok{\#current column names}
\end{Highlighting}
\end{Shaded}

\begin{verbatim}
 [1] "Subcorpus"                                                     
 [2] "Filename"                                                      
 [3] "Year.data.collection"                                          
 [4] "Placement.test.score..raw."                                    
 [5] "Placement.test.score...."                                      
 [6] "Proficiency"                                                   
 [7] "Sex"                                                           
 [8] "Age"                                                           
 [9] "School.University.Institution"                                 
[10] "Major"                                                         
[11] "Year.at.university.school"                                     
[12] "L1"                                                            
[13] "Father.s.native.language"                                      
[14] "Mother.s.nativelanguage"                                       
[15] "Languages.spoken.at.home"                                      
[16] "Age.of.exposure.to.Spanish"                                    
[17] "Years.studying.Spanish"                                        
[18] "Stay.abroad.in.Spanish.speaking.country.....1.month."          
[19] "Stay.abroad..where."                                           
[20] "Stay.abroad..when."                                            
[21] "Stay.abroad..months."                                          
[22] "Language.certificates..type.and.level."                        
[23] "Proficiency..self.assessment..speaking"                        
[24] "Proficiency..self.assessment..listening"                       
[25] "Proficiency..self.assessment..reading"                         
[26] "Proficiency..self.assessment..writing"                         
[27] "Proficiency..self.assessment."                                 
[28] "Additional.foreign.language.s."                                
[29] "Proficiency..self.assessment..in.additional.language.speaking" 
[30] "Proficiency..self.assessment..in.additional.language.listening"
[31] "Proficiency..self.assessment..in.additional.language.reading"  
[32] "Proficiency..self.assessment..in.additional.language.writing"  
[33] "Medium"                                                        
[34] "Task.number"                                                   
[35] "Task.title"                                                    
[36] "Writting.audio.details"                                        
[37] "Minutes.taken.to.complete.the.task"                            
[38] "Where.the.task.was.done"                                       
[39] "Resources.used"                                                
[40] "Text"                                                          
[41] "Original.text"                                                 
\end{verbatim}

\begin{Shaded}
\begin{Highlighting}[]
\NormalTok{englishL1\_data }\OtherTok{\textless{}{-}}\NormalTok{ englishL1\_data }\SpecialCharTok{|\textgreater{}}  \CommentTok{\#renaming columns as in Kochers df}
  \FunctionTok{rename}\NormalTok{(}
    \StringTok{"Year"} \OtherTok{=} \StringTok{"Year.data.collection"}\NormalTok{,}
    \StringTok{"AgeExposureSpanish"} \OtherTok{=} \StringTok{"Age.of.exposure.to.Spanish"}\NormalTok{,}
    \StringTok{"YearsstudySpanish"} \OtherTok{=} \StringTok{"Years.studying.Spanish"}\NormalTok{,}
    \StringTok{"AbroadSpanishCountry"} \OtherTok{=} \StringTok{"Stay.abroad.in.Spanish.speaking.country.....1.month."}\NormalTok{,}
    \StringTok{"DurationTask"} \OtherTok{=} \StringTok{"Minutes.taken.to.complete.the.task"}\NormalTok{,}
    \StringTok{"LocationTask"} \OtherTok{=} \StringTok{"Where.the.task.was.done"}\NormalTok{,}
    \StringTok{"ForeignLanguage"} \OtherTok{=} \StringTok{"Additional.foreign.language.s."}\NormalTok{,}
    \StringTok{"Score"} \OtherTok{=} \StringTok{"Placement.test.score...."}\NormalTok{,}
    \StringTok{"Place"} \OtherTok{=} \StringTok{"School.University.Institution"}\NormalTok{,}
    \StringTok{"FullText"} \OtherTok{=} \StringTok{"Text"}
\NormalTok{  ) }\SpecialCharTok{|\textgreater{}} 
  \FunctionTok{select}\NormalTok{( }\CommentTok{\#deselecting columns not needed}
    \SpecialCharTok{{-}}\NormalTok{Subcorpus, }\SpecialCharTok{{-}}\NormalTok{Placement.test.score..raw., }\SpecialCharTok{{-}}\NormalTok{L1, }\SpecialCharTok{{-}}\NormalTok{Major, }\SpecialCharTok{{-}}\NormalTok{Year.at.university.school, }\SpecialCharTok{{-}}\NormalTok{Father.s.native.language, }\SpecialCharTok{{-}}\NormalTok{Mother.s.nativelanguage, }\SpecialCharTok{{-}}\NormalTok{Stay.abroad..where., }\SpecialCharTok{{-}}\NormalTok{Stay.abroad..when., }\SpecialCharTok{{-}}\NormalTok{Stay.abroad..months., }\SpecialCharTok{{-}}\NormalTok{Language.certificates..type.and.level., }\SpecialCharTok{{-}}\NormalTok{Proficiency..self.assessment..speaking, }\SpecialCharTok{{-}}\NormalTok{Proficiency..self.assessment..listening, }\SpecialCharTok{{-}}\NormalTok{Proficiency..self.assessment..reading, }\SpecialCharTok{{-}}\NormalTok{Proficiency..self.assessment..writing, }\SpecialCharTok{{-}}\NormalTok{Proficiency..self.assessment., }\SpecialCharTok{{-}}\NormalTok{Proficiency..self.assessment..in.additional.language.speaking, }\SpecialCharTok{{-}}\NormalTok{Proficiency..self.assessment..in.additional.language.listening, }\SpecialCharTok{{-}}\NormalTok{Proficiency..self.assessment..in.additional.language.reading, }\SpecialCharTok{{-}}\NormalTok{Proficiency..self.assessment..in.additional.language.writing, }\SpecialCharTok{{-}}\NormalTok{Medium, }\SpecialCharTok{{-}}\NormalTok{Task.number, }\SpecialCharTok{{-}}\NormalTok{Task.title, }\SpecialCharTok{{-}}\NormalTok{Writting.audio.details, }\SpecialCharTok{{-}}\NormalTok{Resources.used, }\SpecialCharTok{{-}}\NormalTok{Original.text) }\SpecialCharTok{|\textgreater{}} 
  \FunctionTok{mutate}\NormalTok{( }\CommentTok{\#adding ProficiencyNr column and values 1{-}6}
    \AttributeTok{ProficiencyNr =} \FunctionTok{case\_when}\NormalTok{(}
\NormalTok{      Proficiency }\SpecialCharTok{==} \StringTok{"Lower beginner"} \SpecialCharTok{\textasciitilde{}} \StringTok{"1"}\NormalTok{,}
\NormalTok{      Proficiency }\SpecialCharTok{==} \StringTok{"Upper beginner"} \SpecialCharTok{\textasciitilde{}} \StringTok{"2"}\NormalTok{,}
\NormalTok{      Proficiency }\SpecialCharTok{==} \StringTok{"Lower intermediate"} \SpecialCharTok{\textasciitilde{}} \StringTok{"3"}\NormalTok{,}
\NormalTok{      Proficiency }\SpecialCharTok{==} \StringTok{"Upper intermediate"}\SpecialCharTok{\textasciitilde{}} \StringTok{"4"}\NormalTok{,}
\NormalTok{      Proficiency }\SpecialCharTok{==} \StringTok{"Lower advanced"} \SpecialCharTok{\textasciitilde{}} \StringTok{"5"}\NormalTok{,}
\NormalTok{      Proficiency }\SpecialCharTok{==} \StringTok{"Upper advanced"} \SpecialCharTok{\textasciitilde{}} \StringTok{"6"}\NormalTok{,}
      \ConstantTok{TRUE} \SpecialCharTok{\textasciitilde{}}\NormalTok{ Proficiency)}
\NormalTok{    )}
\end{Highlighting}
\end{Shaded}

Next, I add the columns `NSLhome' and `NSLadditional':

\begin{Shaded}
\begin{Highlighting}[]
\FunctionTok{table}\NormalTok{(englishL1\_data}\SpecialCharTok{$}\NormalTok{Languages.spoken.at.home) }\CommentTok{\# displaying languages spoken at home}
\end{Highlighting}
\end{Shaded}

\begin{verbatim}

                 English  English, Farsi, Kurdish          English, French 
                     147                        1                        1 
English, Gujarati, Hindi        English, Japanese         English, Spanish 
                       1                        2                        2 
         Hindi, Gujarati         Italian, English               Portuguese 
                       1                        1                        1 
                 Punjabi                  Spanish         Spanish, English 
                       1                        2                        1 
\end{verbatim}

\begin{Shaded}
\begin{Highlighting}[]
\CommentTok{\#Apart from English, some speakers also grew up acquiring Farsi, Kurdish, French, Gujarati, Hindi, Japanese, Spanish, Italian, Portuguese and Punjabi}

\NormalTok{nsl\_home\_L1engl }\OtherTok{\textless{}{-}} \FunctionTok{c}\NormalTok{(}\StringTok{"Farsi"}\NormalTok{, }\StringTok{"Kurdish"}\NormalTok{, }\StringTok{"Gujarati"}\NormalTok{, }\StringTok{"Hindi"}\NormalTok{, }\StringTok{"Japanese"}\NormalTok{, }\StringTok{"Spanish"}\NormalTok{, }\StringTok{"Italian"}\NormalTok{, }\StringTok{"Portuguese"}\NormalTok{, }\StringTok{"Punjabi"}\NormalTok{)}

\FunctionTok{class}\NormalTok{(englishL1\_data}\SpecialCharTok{$}\NormalTok{Languages.spoken.at.home) }\CommentTok{\#vector type "character"}
\end{Highlighting}
\end{Shaded}

\begin{verbatim}
[1] "character"
\end{verbatim}

\begin{Shaded}
\begin{Highlighting}[]
\NormalTok{englishL1\_data }\OtherTok{\textless{}{-}}\NormalTok{ englishL1\_data }\SpecialCharTok{|\textgreater{}}
  \FunctionTok{mutate}\NormalTok{(}
    \AttributeTok{NSLhome =} \FunctionTok{case\_when}\NormalTok{(}
      \FunctionTok{str\_detect}\NormalTok{(Languages.spoken.at.home, }\FunctionTok{paste}\NormalTok{(nsl\_home\_L1engl, }\AttributeTok{collapse =} \StringTok{"|"}\NormalTok{)) }\SpecialCharTok{\textasciitilde{}} \StringTok{"Yes"}\NormalTok{,}
      \ConstantTok{TRUE} \SpecialCharTok{\textasciitilde{}} \StringTok{"No"}
\NormalTok{    )}
\NormalTok{  ) }\CommentTok{\#If the speakers grew up with a NSL at home the answer is "Yes", else "No"}


\CommentTok{\#Adding the NSLadditional column}

\FunctionTok{class}\NormalTok{(englishL1\_data}\SpecialCharTok{$}\NormalTok{ForeignLanguage) }\CommentTok{\#vector type "character"}
\end{Highlighting}
\end{Shaded}

\begin{verbatim}
[1] "character"
\end{verbatim}

\begin{Shaded}
\begin{Highlighting}[]
\FunctionTok{table}\NormalTok{(englishL1\_data}\SpecialCharTok{$}\NormalTok{ForeignLanguage) }\CommentTok{\# displaying foreign languages learned}
\end{Highlighting}
\end{Shaded}

\begin{verbatim}

                                           
                                       123 
                    American Sign Language 
                                         1 
                                    Arabic 
                                         2 
Arabic (Egyptian Colloquial), Arabic (MSA) 
                                         1 
                              Arabic (MSA) 
                                         1 
                                   Catalan 
                                         1 
                                    Creole 
                                         1 
                                    French 
                                         5 
                                    German 
                                         5 
                            German, French 
                                         1 
                           Gujarati, Hindi 
                                         1 
                                     Hindi 
                                         1 
                                   Italian 
                                         2 
                              Italian (A1) 
                                         1 
                           Italian, French 
                                         1 
                                  Japanese 
                                         3 
              Japanese, Portuguese, French 
                                         1 
                                 Kaqchikel 
                                         1 
                                   Kurdish 
                                         1 
                                 Norwegian 
                                         1 
                                Portuguese 
                                         1 
                           Portuguese (A1) 
                                         1 
                                   Russian 
                                         3 
                                   Swahili 
                                         2 
\end{verbatim}

\begin{Shaded}
\begin{Highlighting}[]
\NormalTok{nsl\_foreignlg\_L1engl }\OtherTok{\textless{}{-}} \FunctionTok{c}\NormalTok{(}\StringTok{"Arabic"}\NormalTok{, }\StringTok{"Catalan"}\NormalTok{, }\StringTok{"Gujarati"}\NormalTok{, }\StringTok{"Hindi"}\NormalTok{, }\StringTok{"Italian"}\NormalTok{, }
                   \StringTok{"Japanese"}\NormalTok{, }\StringTok{"Portuguese"}\NormalTok{, }\StringTok{"Kaqchikel"}\NormalTok{, }\StringTok{"Kurdish"}\NormalTok{, }
                   \StringTok{"Russian"}\NormalTok{, }\StringTok{"Swahili"}\NormalTok{)}


\NormalTok{englishL1\_data }\OtherTok{\textless{}{-}}\NormalTok{ englishL1\_data }\SpecialCharTok{|\textgreater{}}
  \FunctionTok{mutate}\NormalTok{(}\AttributeTok{NSLadditional =} \FunctionTok{case\_when}\NormalTok{(}
\NormalTok{    NSLhome }\SpecialCharTok{==} \StringTok{"Yes"} \SpecialCharTok{\textasciitilde{}} \StringTok{"Yes"}\NormalTok{,}
\NormalTok{    NSLhome }\SpecialCharTok{==} \StringTok{"No"} \SpecialCharTok{\&} \FunctionTok{str\_detect}\NormalTok{(ForeignLanguage, }\FunctionTok{paste}\NormalTok{(nsl\_foreignlg\_L1engl, }\AttributeTok{collapse =} \StringTok{"|"}\NormalTok{)) }\SpecialCharTok{\textasciitilde{}} \StringTok{"Yes"}\NormalTok{,}
      \ConstantTok{TRUE} \SpecialCharTok{\textasciitilde{}} \StringTok{"No"}\NormalTok{)}
\NormalTok{    ) }\CommentTok{\#When speakers already acquired a NSL at home than column "NSLadditional" will already be "yes". If they did not but learned a NSL later (in this case Arabic, Catalan, Gujarati, Hindi, Italian, Japanese, Portuguese, Kaqchikel, Kurdish, Russian, Swahili), than in column "NSLadditional" will be annotated "yes".}


\CommentTok{\#Checking if the vector type of the column \textquotesingle{}ProficiencyNr\textquotesingle{} aligns with the L1German dataset: }

\FunctionTok{class}\NormalTok{(englishL1\_data}\SpecialCharTok{$}\NormalTok{ProficiencyNr) }\CommentTok{\#character}
\end{Highlighting}
\end{Shaded}

\begin{verbatim}
[1] "character"
\end{verbatim}

\begin{Shaded}
\begin{Highlighting}[]
\NormalTok{englishL1\_data }\OtherTok{\textless{}{-}}\NormalTok{ englishL1\_data }\SpecialCharTok{|\textgreater{}} 
  \FunctionTok{mutate}\NormalTok{(}\AttributeTok{ProficiencyNr =} \FunctionTok{as.integer}\NormalTok{(ProficiencyNr)) }\CommentTok{\#changed to integer now}

\CommentTok{\#write\_csv(englishL1\_data, here("data", "englishL1\_data.csv")) \#saving added data as .csv file}
\end{Highlighting}
\end{Shaded}

I segmented each clause per row of every text using the machine-based
tool \emph{Meta Llama 3.3 70B} and \emph{DeepSeek R1 Distill LLama 70B}
Instruct for this process via the University of Cologne's academic cloud
access. The englishL1.csv\_segmented df is then read into R again for
merging.

\begin{Shaded}
\begin{Highlighting}[]
\CommentTok{\#Reading the segmented file into R}

\NormalTok{englishL1\_segmented }\OtherTok{\textless{}{-}} \FunctionTok{read.csv}\NormalTok{(here}\SpecialCharTok{::}\FunctionTok{here}\NormalTok{(}\StringTok{"data"}\NormalTok{, }\StringTok{"englishL1\_segmented.csv"}\NormalTok{), }\AttributeTok{stringsAsFactors=}\ConstantTok{TRUE}\NormalTok{)}

\FunctionTok{colnames}\NormalTok{(englishL1\_segmented)}
\end{Highlighting}
\end{Shaded}

\begin{verbatim}
[1] "Text"     "Filename"
\end{verbatim}

\begin{Shaded}
\begin{Highlighting}[]
\FunctionTok{colnames}\NormalTok{(englishL1\_data) }
\end{Highlighting}
\end{Shaded}

\begin{verbatim}
 [1] "Filename"                 "Year"                    
 [3] "Score"                    "Proficiency"             
 [5] "Sex"                      "Age"                     
 [7] "Place"                    "Languages.spoken.at.home"
 [9] "AgeExposureSpanish"       "YearsstudySpanish"       
[11] "AbroadSpanishCountry"     "ForeignLanguage"         
[13] "DurationTask"             "LocationTask"            
[15] "FullText"                 "ProficiencyNr"           
[17] "NSLhome"                  "NSLadditional"           
\end{verbatim}

\begin{Shaded}
\begin{Highlighting}[]
\CommentTok{\# Combining the segmented clauses of the dataframes \textquotesingle{}englishL1\_segmented\textquotesingle{} with the metadata of the dataframes \textquotesingle{}englishL1\_data\textquotesingle{}}

\NormalTok{joined\_englishL1\_data }\OtherTok{\textless{}{-}} \FunctionTok{inner\_join}\NormalTok{(}\AttributeTok{x =}\NormalTok{ englishL1\_segmented, }\AttributeTok{y =}\NormalTok{ englishL1\_data, }\AttributeTok{by =} \StringTok{"Filename"}\NormalTok{)}

\CommentTok{\# Adding the \textquotesingle{}SubjectGenus\textquotesingle{} variable to the \textquotesingle{}joined\_germanL1\_data\textquotesingle{} dataframe for manual annotation}

\NormalTok{joined\_englishL1\_data[, }\StringTok{\textquotesingle{}SubjectGenus\textquotesingle{}}\NormalTok{] }\OtherTok{=} \ConstantTok{NA}
\NormalTok{joined\_englishL1\_data[, }\StringTok{\textquotesingle{}Textnr\textquotesingle{}}\NormalTok{] }\OtherTok{=} \ConstantTok{NA}
\NormalTok{joined\_englishL1\_data[, }\StringTok{\textquotesingle{}Sentence\textquotesingle{}}\NormalTok{] }\OtherTok{=} \ConstantTok{NA}
\NormalTok{joined\_englishL1\_data[, }\StringTok{\textquotesingle{}ID\textquotesingle{}}\NormalTok{] }\OtherTok{=} \ConstantTok{NA}
\NormalTok{joined\_englishL1\_data[, }\StringTok{\textquotesingle{}SubjectPosition\textquotesingle{}}\NormalTok{] }\OtherTok{=} \ConstantTok{NA}
\NormalTok{joined\_englishL1\_data[, }\StringTok{\textquotesingle{}ClauseType\textquotesingle{}}\NormalTok{] }\OtherTok{=} \ConstantTok{NA}
\NormalTok{joined\_englishL1\_data[, }\StringTok{\textquotesingle{}Specificity\textquotesingle{}}\NormalTok{] }\OtherTok{=} \ConstantTok{NA}
\NormalTok{joined\_englishL1\_data[, }\StringTok{\textquotesingle{}InformationStructure\textquotesingle{}}\NormalTok{] }\OtherTok{=} \ConstantTok{NA}
\NormalTok{joined\_englishL1\_data[, }\StringTok{\textquotesingle{}ReferentialContinuity\textquotesingle{}}\NormalTok{] }\OtherTok{=} \ConstantTok{NA}
\NormalTok{joined\_englishL1\_data[, }\StringTok{\textquotesingle{}Subjecttype\textquotesingle{}}\NormalTok{] }\OtherTok{=} \ConstantTok{NA}
\NormalTok{joined\_englishL1\_data[, }\StringTok{\textquotesingle{}Verb\textquotesingle{}}\NormalTok{] }\OtherTok{=} \ConstantTok{NA}
\NormalTok{joined\_englishL1\_data[, }\StringTok{\textquotesingle{}VerbConstruction\textquotesingle{}}\NormalTok{] }\OtherTok{=} \ConstantTok{NA}
\end{Highlighting}
\end{Shaded}

As the L1German data - including the newly added speaker texts - contain
2997 clauses, I aim to extract a similarly high number of L1English
clauses that I annotate manually. Therefore, I select a random sample of
2000 clauses in the following code lines:

\begin{Shaded}
\begin{Highlighting}[]
\NormalTok{joined\_englishL1\_data}\SpecialCharTok{$}\NormalTok{sample }\OtherTok{\textless{}{-}} \ConstantTok{NA} \CommentTok{\#creating sample column}
\FunctionTok{set.seed}\NormalTok{(}\DecValTok{42}\NormalTok{)  }\CommentTok{\#Setting a seed to ensure that the outcome of my randomly generated number series are always exactly the same, so that my study is reproducible }
\NormalTok{random\_clauses }\OtherTok{\textless{}{-}} \FunctionTok{sample}\NormalTok{(}\FunctionTok{nrow}\NormalTok{(joined\_englishL1\_data), }\DecValTok{2000}\NormalTok{)}

\CommentTok{\#Assigning the word \textquotesingle{}sample\textquotesingle{} to the selected rows}

\StringTok{"sample"} \OtherTok{{-}\textgreater{}}\NormalTok{ joined\_englishL1\_data}\SpecialCharTok{$}\NormalTok{sample[random\_clauses]}

\CommentTok{\#write\_csv(joined\_englishL1\_data, here("data", "joined\_englishL1\_data.csv")) \#saving added data as .csv file}
\end{Highlighting}
\end{Shaded}

\textbf{❓PRE-ANNOTATION QUESTIONS❓}

\begin{itemize}
\item
  When L1 English speakers used English words, especially verbs as a
  lack of vocabulary knowledge, should I still annotate as if it was the
  Spanish word (e.g.~\emph{y Charlie gives un hombre él bebe}
  (EN\_WR\_24\_16\_0.5\_14\_SZG))? I haven't discovered this issue in
  the L1 German data so far to be able to compare it with Kocher's
  results.

  -\textgreater{} I will exclude these cases containing English verbs.
\item
  The variable NSLadditional in Kocher's L1 German speaker dataset is
  unclear to me. For example, the speaker `DE\_WR\_27\_19\_6\_14\_AOK'
  grew up speaking Russian and German (which is why NSLhome is set to
  `Yes' for Russian), but apart from Spanish, they only speak English.
  Nevertheless, Kocher assigned the value `Yes' to NSLadditional, even
  though English is not an NSL and there are no other languages studied
  that might be an NSL. Another example: The speaker
  `DE\_WR\_38\_21\_7\_14\_CW' learnt Hebrew alongside Spanish and did
  not mention English at all. Nevertheless, Kocher assigned `English' to
  the ForeignLanguage column and added `Yes' to NSLadditional. I don't
  understand where the languages mentioned in the `ForeignLanguages'
  column or the `NSLadditional' categorisation as `Yes/No' are coming
  from. Another example regarding speaker `DE\_WR\_38\_21\_10\_14\_HYR':
  There, Kocher classified `Hebrew' as `ForeignLanguage' and
  `NSLadditional' as `No', even though the person studied English and
  never mentions Hebrew in the raw data that I extracted. Could the
  metadata be incorrect?

  --\textgreater{} In some cases the data is incorrect. But:
  ``NSLadditional'' is a general variable that is classified as ``yes'',
  when the speaker acquired a NSL at home or studied it afterwards. So
  when the column ``NSLhome'' is set to ``yes'', ``NSLadditional'' will
  always be ``yes'' too.
\item
  Further, I annotated the column `Textnr' not only in the new L1 German
  data but re-annotated it also in Kocher's pre-annotated L1 data (both
  merged in dataframe joined\_germanL1\_data) as it was partially
  incorrect.
\end{itemize}

\textbf{ANNOTATION SCHEME:}

{\def\LTcaptype{none} % do not increment counter
\begin{longtable}[]{@{}
  >{\raggedright\arraybackslash}p{(\linewidth - 2\tabcolsep) * \real{0.0996}}
  >{\raggedright\arraybackslash}p{(\linewidth - 2\tabcolsep) * \real{0.9004}}@{}}
\toprule\noalign{}
\begin{minipage}[b]{\linewidth}\raggedright
Variable
\end{minipage} & \begin{minipage}[b]{\linewidth}\raggedright
Annotated values
\end{minipage} \\
\midrule\noalign{}
\endhead
\bottomrule\noalign{}
\endlastfoot
\textbf{SubjectPosition} & PRE, POST, NA (in nullsub cases) \\
\textbf{ClauseType} & MC, SC \\
\textbf{Specificity} & DEF, INFED, NA (when nullsub is used) \\
\textbf{InformationStructure} & GI(VEN), NE(W), NA (unknown subject
form) \\
\textbf{ReferentialContinuity} & Y(ES), N(O), NA (unknown subject
form) \\
\textbf{SubjectType} & nullsub, dp, pron \\
\textbf{Verb} & Verb of the corresponding sentence \\
\textbf{VerbConstruction} & simple-finite, periphrasis-lexical,
periphrasis-gerund, modal, infinitive \\
\textbf{SubjectGenus} & \begin{minipage}[t]{\linewidth}\raggedright
masc, fem, unknown¹, general

¹Some of the \emph{unknown} cases, that Kocher actually claims to have
excluded, are still within the annotated data. \emph{Unknown} includes
cases in which the subject is not clearly identifiable from the prior
context, where an impersonal form has been used.

\emph{(General} refers to first person plural forms as in \emph{Here we
can see {[}\ldots{]}}.)\\
\textbf{\textgreater{} SOLCHE FÄLLE GANZ RAUS?}

\emph{Unknown} examples:

\begin{itemize}
\item
  \emph{cuando de repente le tiran como basura} from
  kocher\_native\_data 1\_5;
\item
  \emph{Se le ha caído algo} from kocher\_native\_data 2\_10;
\item
  \emph{que han de cuidarlo y darle mimos} from kocher\_native\_data
  7\_25;
\item
  \emph{y alguien vuelve a tirar basura muy cerca de él} from
  kocher\_native\_data 10\_7
\end{itemize}

General\strut
\end{minipage} \\
\end{longtable}
}

\textbf{ANNOTATIONSSCHEMA:}

{\def\LTcaptype{none} % do not increment counter
\begin{longtable}[]{@{}
  >{\raggedright\arraybackslash}p{(\linewidth - 2\tabcolsep) * \real{0.2692}}
  >{\raggedright\arraybackslash}p{(\linewidth - 2\tabcolsep) * \real{0.7308}}@{}}
\toprule\noalign{}
\begin{minipage}[b]{\linewidth}\raggedright
Variable
\end{minipage} & \begin{minipage}[b]{\linewidth}\raggedright
Annotierte Werte
\end{minipage} \\
\midrule\noalign{}
\endhead
\bottomrule\noalign{}
\endlastfoot
\textbf{SubjectPosition} & PRE, POST, NA (bei nullsubj) \\
\textbf{ClauseType} & MC, SC \\
\textbf{Specificity} & DEF, INFED, NA (bei nullsub) \\
\textbf{InformationStructure} & GI(VEN), NE(W), NA (bei unbekannter
Subjektform) \\
\textbf{ReferentialContinuity} & Y(ES), N(O), NA (bei unbekannter
Subjektform) \\
\textbf{SubjectType} & nullsub, dp, pron \\
\textbf{Verb} & Verb des jeweiligen Clauses \\
\textbf{VerbConstruction} & simple-finite, periphrasis-lexical,
periphrasis-gerund, modal, infinitive \\
\textbf{SubjectGenus} & masc, fem, unknown, general \\
\end{longtable}
}

Some notes to explain the variables in Kocher
(\citeproc{ref-kocher2025}{2025b})'s and my data annotation:

\begin{itemize}
\item
  The `Foreign language' column seems to refer only to additional
  foreign languages studied, rather than to other L1s acquired at home
  (NSLhome). Kocher (\citeproc{ref-kocher2025}{2025b}) simply
  categorised whether speakers had acquired a null subject language at
  home, without specifying which one. In this case, `No' could mean that
  the speaker has not acquired another language at all, or that the
  language they grew up with was not a null subject language.
\item
  Kocher (\citeproc{ref-kocher2025}{2025b}) separated the subordinate
  clauses (SCs), such as `Mientras la gente {[}que vive en los
  edificios{]} tiran su basura de la ventana a la calle', from the main
  clause (MC) and displayed them in the row below.
\item
  Cut-out SCs were marked as {[}\ldots{]} within the MCs.
\item
  Direct speech was also marked in brackets and not annotated. Lines
  containing English direct speech were omitted
  (e.g.~kocher\_data\_germanL1 ID9\_28), but not when the direct speech
  was translated into Spanish (e.g.~kocher\_data\_germanL1 ID59\_9). The
  annotation was inconsistent: for example, in line 2328, direct speech
  was marked with `{[}\ldots{]}', as I did too.
\item
  Further, I excluded meta commentaries in the L1 English dataset, such
  as \emph{La pelicula es en negro y blanco, pero es muy beuno} (in
  EN\_WR\_14\_23\_1\_14\_BB), as well as comments written in the first
  person singular or plural, such as \emph{i do not know what else to
  write i already did my best with the previous sentences i wrote}
  (EN\_WR\_20\_21\_4\_14\_PKO) or \emph{Al inizio del short lo vemos
  vestido mal con la ropa desarreglada} (EN\_WR\_39\_62\_2\_14\_MA)
  before the annotation process, in order to avoid their accidental
  inclusion in the random sample. I decided to exclude them, as there
  are too few first-person sentences to obtain reliable results about
  these subjects' behaviour and make meaningful model interpretations.
  In contrast, Kocher (\citeproc{ref-kocher2025}{2025b}) mostly deleted
  them after assigning an ID number, as some ID numbers are missing from
  the final annotated document (see the kocher\_germanL1\_data file).
  However, some clauses containing first-person singular or plural
  subjects remained in her annotation and consequently influenced the
  model (see e.g.~\emph{i creó} in joined\_germanL1\_data\_annotation
  (DE\_WR\_31\_27\_24\_14\_MS no. 19\_2), in
  joined\_germanL1\_data\_annotation (\emph{Nunca aprendì espanol}
  DE\_WR\_20\_50\_2\_14\_HO no. 4\_2).
\item
  I took over the same exclusion criteria as Kocher
  (\citeproc{ref-kocher_subject_2025}{2025a}) which applies to
  imperatives, indefinite ``se'' constructions and impersonal sentences
  (e.g., hay, llueve). Also I decided to exclude sentences that do not
  contain a verb (see \emph{una película sobre un bebé encontrado y
  buscando amor} EN\_WR\_12\_30\_0.7\_14\_NC no. 6\_1) or in which the
  student used an English word (see \emph{La cop chase the el hombre}
  EN\_WR\_12\_30\_0.7\_14\_NC no. 6\_5). These sentences were not
  deleted but taking off the \emph{sample} selection in order to still
  be able to follow the referential continuity and information structure
  of the subsequent sentences.
\end{itemize}

\protect\phantomsection\label{refs}
\begin{CSLReferences}{1}{0}
\bibitem[\citeproctext]{ref-kocher_subject_2025}
Kocher, A. (2025a). The subject realization in {L2} {Spanish} by
{German} {L1} speakers: {A} corpus study. \emph{Isogloss. Open Journal
of Romance Linguistics}, \emph{11}(2), 1--34.
\url{https://doi.org/10.5565/rev/isogloss.437}

\bibitem[\citeproctext]{ref-kocher2025}
Kocher, A. (2025b). The subject realization in L2 spanish by german L1
speakers: A corpus study. \emph{Isogloss. Open Journal of Romance
Linguistics}, \emph{11}(2), 1--34.
\url{https://doi.org/10.5565/rev/isogloss.437}

\end{CSLReferences}




\end{document}
